% Background: technical_background sec 3 (inscribed-sphere radius). The scaling law is defined here,
% not in a background chapter (user, 2026-10-05; moved from clinical_background_v2/04).
% CITATION WANTED: Murray (1926) for the exponent 3; no refs.bib entry yet, so it is stated uncited.
\subsection{Daughter Ratio}
\label{sec:method-ratio}

The radii at a bifurcation were compared with the scaling law followed by healthy branching trees,
according to which the radius of the parent raised to a power equals the sum of the daughter radii
raised to the same power.
Murray derived a power of 3, whereas pooling 1070 coronary trees gave 2.39 (95\% confidence interval
2.24 to 2.54) \cite{taylor2024murray}, close to the 7/3 that Huo and Kassab derived from the minimum
energy a vascular tree requires \cite{huo2012scaling}.
The departure of a bifurcation from the law of Huo and Kassab was measured by the deviation
\begin{equation}
  \Delta_{HK} = \frac{1}{r_0}\left(\frac{r_1^{7/3} + r_2^{7/3}}{2}\right)^{3/7} - 2^{-3/7},
  \label{eq:method-hk}
\end{equation}
where $r_0$ is the parent radius and $r_1$ and $r_2$ are the daughter radii, so that
$\Delta_{HK} = 0$ for a bifurcation that follows the law.
